\chapter{面部表情分类}
\label{chap:fer}

\section{任务与数据划分}
本任务以单张灰度面部图像为输入，在愤怒、恐惧、快乐、悲伤、惊讶和中性
六类中输出一个标签。课程数据包含 28\,275 张带标注训练图和 7\,065 张
测试图，分辨率均为 $48\times48$。官方 FER2013 的原始任务包含七种类别，
包括课程目录中未出现的厌恶类\citep{goodfellow2015fer}。因此，本文仅用
原论文说明数据背景，训练和测试统计均由课程文件夹与样例提交核查。
图~\ref{fig:fer-data} 左侧给出实际样本计数。快乐样本有 7\,192 张，惊讶
样本有 3\,202 张，两者相差约 2.25 倍。整体准确率会更多反映高频类的表现，
故错误分析另外报告逐类召回率与宏平均 F1。

数据层从训练目录建立固定的类索引，依次为 Angry、Fear、Happy、Sad、
Surprise 和 Neutral。Kaggle 提交将其映射为整数 0 至 5，写入
\code{ID,Emotion}。训练集统一使用分层五折，数据划分函数默认种子 42，
分类训练配置中的种子 0 控制运行随机性。各骨干以相同的验证索引汇总 OOF
概率，便于在同一批样本上比较成员互补性。

\chart{fer_data.pdf}{训练集类别计数与代表性线上提交}{fig:fer-data}
{左图按竞赛标签编码排序。右图为同一赛题的公开榜记录。图轴放大以显示候选间差异。}

\section{模型与计算流程}
\subsection{输入与迁移学习}
将 JPEG 解码为单通道 \code{uint8} 后，脚本对训练及测试张量做缓存，并将
批处理使用的图像张量放在 GPU 上。35\,340 张 $48\times48$ 图像的像素张量
占约 81.4 MB（十进制字节计），此数值由样本量与每像素 1 字节计算，未含
索引和训练激活。训练时在 GPU 上执行批量仿射、镜像、亮度及对比度变化和
随机擦除。对于 ImageNet 骨干，将灰度张量扩展为三通道并插值到所需尺寸。
扩大输入允许沿用预训练骨干的分辨率配方，并改变早期特征图的采样格点。
原图信息仍受 48 像素分辨率限制。

\archdiagram{fer_arch}{面部表情分类的训练与多模型推理流程}{fig:fer-arch}

ConvNeXt-Tiny 和 ConvNeXt-Small 是主要性能来源。另训练 EfficientNetV2-S、
ResNet 和本地实现的 ResEmoteNet 以比较容量与预测互补性。ResEmoteNet
分支从单通道直接学习，三个卷积块后接带挤压激励的残差块、全局平均池化和
六类线性头\citep{roy2024resemote}。该分支从零训练，保留原生 48 像素输入，
其训练信号与三通道迁移学习成员有所不同。其单模型成绩较低，仍在个别集成中
有非零权重。集成权重只表明它在该联合选择过程中被保留。

\subsection{损失、权重平均与推理}
程序使用平滑交叉熵。Mixup 或 CutMix 时，按每张图的实际混合系数分别
计算两个平滑标签的损失：
\[
\mathcal L=\lambda\,\mathrm{CE}_\varepsilon(\mathbf z,y_a)
 +(1-\lambda)\,\mathrm{CE}_\varepsilon(\mathbf z,y_b),\qquad
 \varepsilon=0.1.
\]
AdamW 对骨干和分类头分别设置学习率。训练更新指数移动平均
$\theta_t^{\rm EMA}=0.999\,\theta_{t-1}^{\rm EMA}+0.001\,\theta_t$。
移动平均减小单批更新对验证模型的波动。脚本会比较即时权重、EMA 和
启用时的快照平均候选。保存最优权重时一并记下来源。预测阶段平均水平
镜像和三个尺度的概率，默认尺度为 0.9、1.0 和 1.1。不同折的概率再按
骨干和输入尺寸分组。

\subsection{融合策略}
对 $M$ 个完整五折模型组，令 $\mathbf p_m(x)$ 是组级概率，
融合概率为 $\sum_{m=1}^M w_m\mathbf p_m(x)$，其中 $w_m\ge 0$ 且
$\sum_m w_m=1$。贪心前向选择逐步加入能提高折外准确率的成员。
概率空间 bagging 多次抽取约 60\% 的折外样本并重复选择，最后平均各次
权重，以减弱个别难例对模型选择的影响。另有 logit 空间的试验候选，
变换方式与概率空间不同。由于折外样本同时参与选择，择优后的 OOF
应作为开发集结果解读。

\section{训练过程与分辨率对照}
\subsection{代表性训练日志}
表~\ref{tab:fer-cfg} 使用 \code{cns224_f0} 运行的配置与日志。
记录覆盖从索引 0 到 47 的 48 轮。训练用 AMP，骨干预热后三个阶段的
关键行为分别是前期学习率爬升、余弦退火和结束前关闭混合增强。
分类头较高的步长使新初始化的线性分类器尽快适配六类数据。验证日志与
\code{summary.json} 对应，EMA 权重在第 12 轮（索引 11）达到 0.73015
并被保存。末轮日志的 EMA 验证准确率为 0.71706，恢复最优权重后的最终
验证准确率为 0.73015。

\begin{table}[H]
\centering\small
\caption{ConvNeXt-Small 224 像素配置第 0 折}\label{tab:fer-cfg}
\begin{tabularx}{\linewidth}{@{}P{48mm}Y@{}}
\toprule
项目 & 配置或记录 \\
\midrule
输入与批大小 & $3\times224\times224$，每批 32 张 \\
训练轮次 & 计划 48 轮，日志 48 轮 \\
学习率与权重衰减 & 骨干 $10^{-4}$，分类头 $10^{-3}$，衰减 0.05 \\
调度与平均 & 3 轮线性预热、余弦退火、EMA 0.999 \\
标签与增强 & 标签平滑 0.1，混合增强概率 0.5，最后 10 轮关闭混合增强 \\
权重选择 & EMA，第 12 轮，验证准确率 0.73015 \\
规模及耗时 & 49.459 M 参数，该折记录 113.17 min \\
\bottomrule
\end{tabularx}
\end{table}

\chart{fer_training.pdf}{ConvNeXt-Small 224 像素第 0 折的训练轨迹}{fig:fer-curve}
{训练准确率按脚本记录，混合增强期间与验证输入的条件不同。EMA 曲线和普通权重曲线分别读取各自的日志字段。}

图~\ref{fig:fer-curve} 中训练损失下降到约 0.433，末轮训练准确率约
0.995，验证准确率保持在约 0.72。训练准确率与验证准确率间的差距、验证峰值
早于末轮、以及末轮 EMA 回落共同表明继续迭代的训练拟合收益较难传到验证集。
因此训练记录应同时给出末轮值与所选权重的值。训练中启用混合增强时的准确率
用混合权重计算，图中对应不同的训练难度，适合观察趋势。

\subsection{分辨率、成本与增强}
原分辨率对照使用同一 ConvNeXt-Tiny 骨干、第 0 折及相同的样本划分。
表~\ref{tab:fer-res} 来自汇总文件。112 升至 224 的验证提升是 0.60 个
百分点，而耗时由 13.89 增至 49.64 分钟，约 3.57 倍。
64 像素 60 轮版本还改变训练轮数，作为低分辨率下的训练预算参照。
在同样的 112 像素第 0 折实验中，几何增强候选的准确率约为
0.7096 至 0.7119，低于原基线 0.7144。局部眉眼和口部结构对仿射变换较
敏感，但单折记录不足以确定是何种扰动引起了下降。

\begin{table}[H]
\centering\small
\caption{ConvNeXt-Tiny 单折分辨率实验}\label{tab:fer-res}
\begin{tabularx}{\linewidth}{@{}Yrrr@{}}
\toprule
运行标记 & 输入边长 & 保存权重准确率 & 记录耗时（min） \\
\midrule
\code{abl_res64} & 64 & 0.6932 & 5.4 \\
\code{abl_res64_e60} & 64 & 0.6999 & 10.7 \\
\code{abl_res112} & 112 & 0.7144 & 13.9 \\
\code{abl_res224} & 224 & 0.7204 & 49.6 \\
\bottomrule
\end{tabularx}
\end{table}

\section{结果与误差分析}
\subsection{单组与集成结果}
五折 ConvNeXt-Small 224 组 OOF 为 0.7260，Small 112 组为 0.7177，
Tiny 224 组为 0.7215。增加较大分辨率 Small 的 \code{v2} 以 0.7372
的 OOF 达到 0.7470 的线上准确率，较早期 \code{v1} 的 0.7430 高
0.40 个百分点。概率空间 bagging 提高 OOF 至 0.7393，线上仍为
0.7470。logit 变体的 OOF 再升至 0.7397，线上为 0.7460。
表~\ref{tab:fer-results} 保留两种并列版本。

\begin{table}[H]
\centering\small
\caption{表情识别的折外准确率与 Kaggle 提交成绩}\label{tab:fer-results}
\begin{tabularx}{\linewidth}{@{}Yrrr@{}}
\toprule
融合版本 & OOF & Public & Private \\
\midrule
v1，112 像素组合 & 0.7278 & 0.7430 & 0.7430 \\
v2，加入 Small 224 & 0.7372 & \textbf{0.7470} & \textbf{0.7470} \\
概率空间 bagging & 0.7393 & \textbf{0.7470} & \textbf{0.7470} \\
logit 空间 bagging & 0.7397 & 0.7460 & 0.7460 \\
原贪心 & 0.7376 & 0.7450 & 0.7450 \\
\bottomrule
\end{tabularx}
\end{table}

最高 OOF 和最高线上分数来自不同的融合配置。权重选择使用同一批
折外概率，变化量也只有千分位量级。结果表明继续优化融合规则的验证增益
难以稳定迁移。两个榜单分数同为 0.7470 也不足以证明两个模型产生了
相同的错误。比较配方时同时考虑了第 0 折耗时与多个骨干的运行成本。

\clearpage
\subsection{类别级误差}
旧贪心汇总的宏平均 F1 为 0.7227。快乐的召回率 0.9110，恐惧为
0.5538，悲伤为 0.6164。快乐的样本数较大且嘴部形态变化明显。恐惧和
惊讶均可能出现睁眼或张嘴，低分辨率与遮挡会弱化判断所需的面部细节。
这里将样本不均衡和局部纹理不足作为可能的原因，逐类汇总只能确证各类表现
不同。\ifLocalReport 实际验证案例见图~\ref{fig:fer-case}。该图来自 \code{cns224_f0}
最优权重的验证概率，服务于错误诊断。表中的集成成绩采用各自独立的
提交记录。\fi
\privatechart{fer_case.png}{ConvNeXt-Small 第 0 折的验证样例}{fig:fer-case}
{每幅图标出真实类别、最大概率类别及该概率。样例用于说明模型即使置信度高也可能分类错误。}

\section{本章小结}
表情识别的可复查改进包括分辨率与耗时对照、同折多骨干概率融合，以及保留 EMA
早期峰值。两种集成在 Kaggle 均达到 0.7470。更复杂的 logit 选择提高了
同批 OOF，但未取得同方向的榜单变化。类别级误差显示，小样本和细粒度表情的
有效识别仍是提升空间。
